\section{Gemini 2.5 Pro Results}
\label{sec:gemini}

\gemini was excluded from the main analysis because it uses approximately 96\% of its completion tokens for internal reasoning (980 of 1{,}020 tokens), leaving only ${\sim}33$ words for the visible response.
This architectural constraint produces severely truncated outputs regardless of topic, creating a confound where brevity scores reflect token allocation rather than avoidance behavior.

For completeness, \tabref{tab:gemini} reports \gemini's raw avoidance scores.
The elevated safe control score ($4.50$) --- compared to $0.05$--$0.40$ for other models --- confirms that brevity alone inflates avoidance scores.
The relative ordering of categories (gray $>$ safe) is consistent with the other models, but we cannot reliably attribute this to avoidance rather than differential truncation effects.

\begin{table}[h]
\centering
\caption{\gemini overall avoidance scores (mean $\pm$ SD). High baseline scores on safe controls are driven by truncation, not avoidance behavior.}
\label{tab:gemini}
\begin{tabular}{@{}lc@{}}
\toprule
\textbf{Category} & \textbf{Avoidance Score} \\
\midrule
Safe Control & $4.50 \pm 1.70$ \\
Gray: Epistemic & $5.20 \pm 1.47$ \\
Gray: Institutional & $5.70 \pm 1.38$ \\
Gray: Social Taboo & $5.80 \pm 1.15$ \\
Known Safety & $9.10 \pm 0.31$ \\
\bottomrule
\end{tabular}
\end{table}

This result highlights an important methodological consideration: behavioral probing via output text is unreliable for reasoning models that allocate most of their token budget to internal chain-of-thought.
Future work should either increase the maximum token budget substantially (e.g., 4{,}096 tokens) or develop probing methods that analyze the reasoning trace itself.
